\documentclass[a4paper, 10pt]{article}
\usepackage{scrextend}
\changefontsizes{9pt}
\input{../course-config}
\usepackage{../vendor/template/CEDT-Assignment-style}

\begin{document}
\subject
\asgntitle{5}{}{Week 5}{\StudentID\ \StudentName}{\DefaultCollaborators}

% ================================================================================ %
%                                    Problem 01                                    %
% ================================================================================ %
\begin{problem}
Generate (7,4) Cyclic code from the data (0101, 1001, 0010, 1110, 0000), and its hardware implementation.
\end{problem}

\begin{solution}
	Take the standard $(7,4)$ Hamming cyclic code, generated by
	$g(x) = x^3 + x + 1$ (binary $1011$), so $n=7$, $k=4$, $n-k=3$. For a
	message $m_3m_2m_1m_0$ write $m(x) = m_3x^3+m_2x^2+m_1x+m_0$. Systematic
	encoding shifts the message into the high-order bits and fills the low
	$3$ bits with a parity check:
	\[
		c(x) = x^{n-k}m(x) + r(x), \qquad r(x) = x^{n-k}m(x) \bmod g(x),
	\]
	so the codeword is $[\,m_3m_2m_1m_0 \mid r_2r_1r_0\,]$, and division is
	mod-2 (XOR instead of subtract/borrow).

	\myspace
	\textbf{Worked example (0101).} $m(x) = x^2+1$, so $x^3m(x) = x^5+x^3$.
	Dividing by $g(x)=x^3+x+1$:
	\begin{align*}
		x^5+x^3            & = \underbrace{(x^2)}_{\text{quotient}}(x^3+x+1) + \underbrace{x^2}_{\text{remainder}}                \\
		(x^2)(x^3+x+1)      & = x^5+x^3+x^2                                                                                       \\
		(x^5+x^3)\ \mathrm{XOR}\ (x^5+x^3+x^2) & = x^2 = r(x)
	\end{align*}
	$\deg(x^2) < \deg g(x)$, so the division stops: $r = 100$, giving
	codeword $c = 0101\,100$. The other four messages divide the same way
	(each verified independently against the systematic generator matrix
	built from $x^{n-k}x^i \bmod g(x)$ for $i=0,1,2,3$):

	\begin{simpletable}{cccc}{$(7,4)$ cyclic codewords, $g(x)=x^3+x+1$}{cyclic-codewords}
		Message $m_3m_2m_1m_0$ & Shifted $x^3m(x)$ & Remainder $r_2r_1r_0$ & Codeword $c$ \\
		\midrule
		0101                    & 0101000            & 100                    & \textbf{0101100} \\
		1001                    & 1001000            & 110                    & \textbf{1001110} \\
		0010                    & 0010000            & 110                    & \textbf{0010110} \\
		1110                    & 1110000            & 100                    & \textbf{1110100} \\
		0000                    & 0000000            & 000                    & \textbf{0000000} \\
	\end{simpletable}

	\myspace
	\textbf{Hardware implementation.} A $(n,k)$ cyclic encoder is built from
	an $(n-k)$-stage feedback shift register whose taps are the coefficients
	$g_0,g_1,g_2$ of $g(x)$ (here $1,1,0$; $g_3=1$ is the implicit feedback
	stage). Let $b_0,b_1,b_2$ be the register stages and $e = G\cdot(m \oplus
	b_2)$ the gated feedback signal, where $G$ is a control line that is $1$
	for the first $k=4$ clocks and $0$ for the last $n-k=3$:
	\[
		b_0 \leftarrow e, \qquad b_1 \leftarrow b_0 \oplus e, \qquad b_2 \leftarrow b_1.
	\]
	While $G=1$ the message bits shift straight through to the output
	\emph{and} into the register (computing the division above); once all
	$4$ message bits have passed, $G$ drops to $0$, which forces $e=0$ (the
	register now just shifts, no more feedback) and switches the output
	to $b_2$, so $b_2,b_1,b_0$ -- exactly $r_2r_1r_0$ -- shift out over the
	next $3$ clocks.

	\begin{figure}[H]
		\centering
		\resizebox{\linewidth}{!}{%
		\begin{tikzpicture}[
				>=Stealth, thick,
				ff/.style={draw, rectangle, minimum width=1.3cm, minimum height=1.1cm, font=\normalsize},
				xorgate/.style={draw, circle, minimum size=0.65cm, font=\normalsize, inner sep=0pt},
				gatebox/.style={draw, rectangle, minimum width=0.8cm, minimum height=0.65cm, font=\small},
				muxbox/.style={draw, rectangle, minimum width=1.0cm, minimum height=2.2cm, font=\small},
				lbl/.style={font=\small},
			]
			% --- main data / division path ---
			\node (min) at (0,0) {$m$};
			\node[xorgate] (x0) at (2.0,0) {$\oplus$};
			\node[gatebox] (g) at (3.8,0) {$G$};
			\coordinate (ebranch) at (5.2,0);
			\node[ff] (ff0) at (6.6,0) {$b_0$};
			\node[xorgate] (x1) at (8.6,0) {$\oplus$};
			\node[ff] (ff1) at (10.6,0) {$b_1$};
			\node[ff] (ff2) at (12.6,0) {$b_2$};

			\draw[->] (min) -- (x0);
			\draw[->] (x0) -- (g);
			\draw (g) -- (ebranch);
			\filldraw (ebranch) circle (1.2pt);
			\draw[->] (ebranch) -- node[lbl, above=1pt]{$e$} (ff0);
			\draw[->] (ff0) -- (x1);
			\draw[->] (x1) -- (ff1);
			\draw[->] (ff1) -- (ff2);

			% feedback tap into the g1 XOR (below the main row)
			\coordinate (edrop) at (5.2,-1.3);
			\coordinate (ecorner) at (8.6,-1.3);
			\draw (ebranch) -- (edrop);
			\draw (edrop) -- (ecorner);
			\draw[->] (ecorner) -- (x1);

			% feedback line from b2 back to the input XOR
			\coordinate (fbtr) at (12.6,2.0);
			\coordinate (fbtl) at (2.0,2.0);
			\draw (ff2) -- (fbtr);
			\draw (fbtr) -- (fbtl);
			\draw[->] (fbtl) -- (x0);

			% --- output selection: MUX picks m (message phase) or b2 (parity phase) ---
			\node[muxbox] (mux) at (12.6,-5.0) {MUX};
			\coordinate (muxb2) at ($(mux.west)+(0,0.65)$);
			\coordinate (muxm) at ($(mux.west)+(0,-0.65)$);

			\draw (ff2) -- (12.6,-4.35);
			\draw[->] (12.6,-4.35) -- (muxb2);
			\node[lbl, right] at (12.6,-2.3) {$b_2$ (parity phase)};

			\draw (min) -- (0,-5.65);
			\draw[->] (0,-5.65) -- (muxm);
			\node[lbl, right] at (0,-3.6) {$m$ (message phase)};

			\draw (g) -- (3.8,-6.1);
			\node[lbl, below] at (3.8,-6.1) {$G$ (select)};
			\draw[->] (3.8,-6.1) -- (mux.south);

			\draw[->] (mux.east) -- ++(1.3,0) node[right, lbl]{$c$};
		\end{tikzpicture}%
		}
		\caption{$3$-stage LFSR encoder for $g(x)=x^3+x+1$: taps at $b_0$
			($g_0{=}1$) and before $b_1$ ($g_1{=}1$), none before $b_2$
			($g_2{=}0$). $G=1$ divides the incoming message; $G=0$ shifts
			the remainder $b_2b_1b_0$ out through the MUX.}
		\label{fig:cyclic-encoder}
	\end{figure}
\end{solution}
% ================================================================================ %

\end{document}
